\setchapterpreamble[u]{\margintoc}
\chapter{Orthogonality, Projection, and Least Squares}
\labch{orthogonality}

\sources{Strang, \emph{Introduction to Linear Algebra}, 6th ed.\ \cite{strang2023ila},
§4.1--4.4; MIT 18.06 \cite{ocw1806}, Lectures~14--17.}

Most systems that arise from measurements have no solution: there are more
equations than unknowns and the data are noisy. Rather than give up, one asks
for the vector that comes closest. Orthogonality is what makes ``closest''
computable.

\section{Inner products, length, and angle}
\labsec{inner-products}

For \(\vect{u},\vect{v}\in\R^{n}\) the inner product is
\(\vect{u}\tp\vect{v}=\sum_i u_iv_i\), the length is
\(\lVert\vect{v}\rVert=\sqrt{\vect{v}\tp\vect{v}}\), and the two vectors are
\emph{orthogonal} when \(\vect{u}\tp\vect{v}=0\).

\begin{proposition}
\(\vect{u}\tp\vect{v}=0\) if and only if
\(\lVert\vect{u}+\vect{v}\rVert^{2}=\lVert\vect{u}\rVert^{2}+\lVert\vect{v}\rVert^{2}\).
\end{proposition}

\begin{proof}
Expanding,
\(\lVert\vect{u}+\vect{v}\rVert^{2}
=(\vect{u}+\vect{v})\tp(\vect{u}+\vect{v})
=\lVert\vect{u}\rVert^{2}+2\vect{u}\tp\vect{v}+\lVert\vect{v}\rVert^{2}\).
The two sides agree exactly when the cross term vanishes.
\end{proof}

\marginnote{So Pythagoras is not an analogy imported from geometry; it is the
algebraic identity above, and it holds in \(\R^{n}\) for every \(n\).}

\section{Orthogonal subspaces}
\labsec{orthogonal-subspaces}

Two subspaces \(V,W\) are orthogonal when \(\vect{v}\tp\vect{w}=0\) for every
\(\vect{v}\in V\) and \(\vect{w}\in W\). The four subspaces of
\refch{four-subspaces} pair up in exactly this way.

\begin{theorem}
For any \(A\in\R^{m\times n}\), the nullspace and the row space are orthogonal
complements in \(\R^{n}\), and the left nullspace and the column space are
orthogonal complements in \(\R^{m}\).
\end{theorem}

\begin{proof}
If \(A\vect{x}=\vect{0}\) then every row of \(A\) has zero inner product with
\(\vect{x}\), hence so does every combination of rows: \(\Null(A)\) is orthogonal
to \(\Col(A\tp)\). Their dimensions are \(n-r\) and \(r\), which sum to \(n\), so
neither can be enlarged inside the other's complement. Applying the same
argument to \(A\tp\) gives the second pair.
\end{proof}

\begin{remark}
This upgrades rank--nullity from an accounting identity to a geometric
statement. The input space splits into a part the matrix reproduces faithfully
and a part it annihilates, and the two parts meet at right angles.
\end{remark}

\section{Projection onto a line}
\labsec{projection-line}

Given \(\vect{a}\neq\vect{0}\), the point of the line \(\{c\vect{a}\}\) closest
to \(\vect{b}\) is the one for which the error \(\vect{b}-c\vect{a}\) is
perpendicular to \(\vect{a}\):
\[
	\vect{a}\tp(\vect{b}-c\vect{a})=0
	\quad\Longrightarrow\quad
	c=\frac{\vect{a}\tp\vect{b}}{\vect{a}\tp\vect{a}},
	\qquad
	\vect{p}=\frac{\vect{a}\tp\vect{b}}{\vect{a}\tp\vect{a}}\,\vect{a}.
\]

Writing this as a matrix acting on \(\vect{b}\) gives the projection matrix
\[
	P=\frac{\vect{a}\vect{a}\tp}{\vect{a}\tp\vect{a}},
	\qquad
	P^{2}=P,\qquad P\tp=P .
\]

\marginnote{\(P^{2}=P\) says projecting twice is projecting once, which any
sensible notion of ``closest point'' must satisfy. Symmetry is the algebraic
trace of the right angle.}

\begin{example}
With \(\vect{a}=[1,1,1]\tp\) and \(\vect{b}=[1,2,2]\tp\),
\(\vect{a}\tp\vect{b}=5\) and \(\vect{a}\tp\vect{a}=3\), so
\(\vect{p}=\tfrac{5}{3}[1,1,1]\tp\) and the error
\(\vect{b}-\vect{p}=\tfrac{1}{3}[-2,1,1]\tp\) indeed satisfies
\(\vect{a}\tp(\vect{b}-\vect{p})=0\).
\end{example}

\section{Projection onto a subspace}
\labsec{projection-subspace}

Replace the single vector by the columns of \(A\in\R^{m\times n}\) and look for
\(\hat{\vect{x}}\) making \(\vect{b}-A\hat{\vect{x}}\) orthogonal to every
column. That is \(A\tp(\vect{b}-A\hat{\vect{x}})=\vect{0}\), or

\begin{equation}
	A\tp A\hat{\vect{x}}=A\tp\vect{b},
	\labeq{normal-equations}
\end{equation}

the \emph{normal equations}. When the columns of \(A\) are independent,
\(A\tp A\) is invertible and
\[
	\hat{\vect{x}}=(A\tp A)^{-1}A\tp\vect{b},
	\qquad
	\vect{p}=A\hat{\vect{x}}=A(A\tp A)^{-1}A\tp\vect{b}.
\]

\begin{proposition}
If the columns of \(A\) are linearly independent then \(A\tp A\) is invertible.
\end{proposition}

\begin{proof}
Suppose \(A\tp A\vect{x}=\vect{0}\). Then
\(\vect{x}\tp A\tp A\vect{x}=\lVert A\vect{x}\rVert^{2}=0\), so
\(A\vect{x}=\vect{0}\), and independence of the columns forces
\(\vect{x}=\vect{0}\). A square matrix with trivial nullspace is invertible.
\end{proof}

\marginnote{The proof is the reason \(A\tp A\) appears everywhere. It converts a
rectangular problem into a square one without losing information, provided the
columns were independent to begin with.}

\section{Least squares}
\labsec{least-squares}

Solving \refeq{normal-equations} minimises \(\lVert A\vect{x}-\vect{b}\rVert^{2}\)
over all \(\vect{x}\). This is the standard way to fit a model with more
observations than parameters.

\begin{example}
Fit a straight line \(b=C+Dt\) through the points \((1,1)\), \((2,2)\),
\((3,2)\). Then
\[
	A=\begin{bmatrix}1&1\\1&2\\1&3\end{bmatrix},
	\qquad
	\vect{b}=\begin{bmatrix}1\\2\\2\end{bmatrix},
	\qquad
	A\tp A=\begin{bmatrix}3&6\\6&14\end{bmatrix},
	\qquad
	A\tp\vect{b}=\begin{bmatrix}5\\11\end{bmatrix}.
\]
The normal equations \(3C+6D=5\) and \(6C+14D=11\) give \(D=\tfrac{1}{2}\) and
\(C=\tfrac{2}{3}\), so the best line is \(b=\tfrac{2}{3}+\tfrac{1}{2}t\). The
residuals are \(-\tfrac{1}{6},\tfrac{1}{3},-\tfrac{1}{6}\); they sum to zero and
satisfy \(A\tp\vect{e}=\vect{0}\), confirming that the error is orthogonal to the
column space.
\end{example}

\begin{remark}
Least squares is the first place where a linear algebra identity and a
minimisation problem coincide. The same coincidence returns in
\refch{optimization} for quadratic objectives and again in \refch{estimation},
where minimising squared error turns out to be maximum likelihood under Gaussian
noise.
\end{remark}

\section{Orthonormal bases and Gram--Schmidt}
\labsec{gram-schmidt}

A set \(\{\vect{q}_1,\ldots,\vect{q}_n\}\) is \emph{orthonormal} when
\(\vect{q}_i\tp\vect{q}_j\) equals \(1\) for \(i=j\) and \(0\) otherwise.
Collecting them as columns of \(Q\) gives \(Q\tp Q=I\), and when \(Q\) is square
this is \(Q^{-1}=Q\tp\).

Orthonormal columns simplify everything above: the normal equations collapse to
\(\hat{\vect{x}}=Q\tp\vect{b}\), and the projection matrix becomes
\(P=QQ\tp\), with no inverse to compute.

\marginnote{Numerically this is the difference between a stable computation and
an unstable one. Forming \(A\tp A\) squares the condition number of \(A\);
orthogonalising first avoids that.}

Gram--Schmidt manufactures an orthonormal basis from any independent set by
subtracting, from each vector in turn, its projection onto everything already
accepted.

\begin{example}
For \(A=\begin{bmatrix}1&1\\1&0\\0&1\end{bmatrix}\), start with
\(\vect{q}_1=\tfrac{1}{\sqrt2}[1,1,0]\tp\). Subtracting the projection of the
second column,
\[
	\begin{bmatrix}1\\0\\1\end{bmatrix}
	-\tfrac{1}{2}\begin{bmatrix}1\\1\\0\end{bmatrix}
	=\begin{bmatrix}\tfrac12\\-\tfrac12\\1\end{bmatrix},
	\qquad
	\vect{q}_2=\tfrac{1}{\sqrt6}\begin{bmatrix}1\\-1\\2\end{bmatrix}.
\]
\end{example}

\section{The factorisation \texorpdfstring{\(A=QR\)}{A=QR}}
\labsec{qr}

Recording the coefficients used during Gram--Schmidt gives a second
factorisation alongside \(A=LU\).

\begin{theorem}
If \(A\in\R^{m\times n}\) has independent columns then \(A=QR\), where the
columns of \(Q\in\R^{m\times n}\) are orthonormal and \(R\in\R^{n\times n}\) is
upper triangular with positive diagonal entries.
\end{theorem}

For the example above,
\[
	R=\begin{bmatrix}\sqrt2 & \tfrac{1}{\sqrt2}\\[2pt] 0 & \sqrt{\tfrac32}\end{bmatrix},
\]
and multiplying \(QR\) returns the two original columns. Substituting \(A=QR\)
into the normal equations turns them into the triangular system
\(R\hat{\vect{x}}=Q\tp\vect{b}\), which is how least squares is solved in
practice.

\begin{remark}
Two factorisations, two purposes. \(A=LU\) records elimination and answers
questions about solving; \(A=QR\) records orthogonalisation and answers
questions about approximating. \refch{svd} produces a third that does both and
survives rank deficiency as well.
\end{remark}
